\par\noindent\textbf{(1)} 一阶Born振幅 f\_B(q)=−[m/(2πℏ²)]∫e\textasciicircum{}\{−i q·r\}V(r)d³r。球对称Fourier积分 ∫e\textasciicircum{}\{−iq·r\}e\textasciicircum{}\{−κr\}/r d³r=4π/(q²+κ²)，因此 f\_B(q)=2mg/[ℏ²(q²+κ²)]，吸引势只改变振幅整体号而对一阶截面无影响。
\par\noindent\textbf{(2)} 微分截面为 dσ/dΩ=|f\_B|²=4m²g²/[ℏ⁴(4k²sin²(θ/2)+κ²)²]。筛选尺度1/κ削弱长程前向峰，θ≈0时截面有限为4m²g²/(ℏ⁴κ⁴)；κ小则前向散射增强。
\par\noindent\textbf{(3)} 积分 dΩ=2π sinθdθ，令u=sin²(θ/2)使sinθdθ=2du，得到 σ\_tot=16πm²g²/ℏ⁴ ∫₀¹du/(4k²u+κ²)²=16πm²g²/[ℏ⁴ κ²(4k²+κ²)]。κ→0时发散，表示无屏蔽Coulomb长程散射总截面不收敛。
\par\noindent\textbf{(4)} Born近似要求散射势相对入射自由传播弱，典型无量纲mg/(ℏ²κ)≪1可保证宽区间有效；当κ→0此条件失效，正是Coulomb势须改用非微扰散射波的原因。低能或近强束缚态阈值时，即使形式截面可算，忽略多次散射仍可能严重错误。
